% Independent preview; original geometry and numeric relationships retained.
% 同一行词元提供查询、键和值；格条展示坐标，权重表按查询位置排列。
\begin{tikzpicture}[x=1mm,y=1mm,font=\figurefont\footnotesize,text=NNInk,
  vector/.style={draw=NNInput,fill=NNInput!18,line width=1.1pt,
    minimum width=28mm,minimum height=7mm,inner sep=0pt},
  read/.style={nn flow,draw=NNHidden,rounded corners=1.2mm},
  mutual/.style={<->,>=Stealth,draw=NNLine!65,line width=.85pt,
    rounded corners=1.2mm},
  auxiliary/.style={nn flow,draw=NNLine!65,line width=.85pt,
    rounded corners=1.2mm}]
  % 一行词元；表示向量在对应词元的正上方。
  \foreach \x/\i/\word in {18/1/猫,55/2/追,92/3/狗}{
    \node[vector] (h\i) at (\x,0) {};
    \foreach \dx in {-7,0,7}{
      \draw[NNInput,line width=.5pt] ({\x+\dx},-3.5)--({\x+\dx},3.5);
    }
    \foreach \j in {1,2,3,4}{
      \node[inner sep=0pt] at ({\x-17.5+7*\j},0) {$h_{\i,\j}$};
    }
    \node[nn heading] (t\i) at (\x,-14) {\word};
    \draw[nn flow,draw=NNWire,line width=.85pt] (t\i.north)--(h\i.south);
  }
  % 双向细线只表示两个位置均可向对方读取；并不表示权重相等。
  \draw[mutual] (h1.east)--(h2.west);
  \draw[mutual] (h2.east)--(h3.west);
  \draw[mutual] (10,3.5)--(10,21)--(100,21)--(100,3.5);
  \draw[auxiliary] (h1.west)--(0,0)--(0,-7)--(10,-7)--(10,-3.5);
  \draw[auxiliary] (h3.east)--(110,0)--(110,-7)--(100,-7)--(100,-3.5);

  % “追”的读取单独展开：每根紫色箭头从值的来源指向其查询位置。
  \draw[read] (h1.north)--(18,10)--(47,10)--(47,3.5);
  \draw[read] (h3.north)--(92,10)--(63,10)--(63,3.5);
  \node[text=NNHidden] at (32.5,14) {$0.25$};
  \node[text=NNHidden] at (77.5,14) {$0.25$};
  \draw[read] (61,-3.5)--(61,-8)--(69,-8)--(69,-3.5);
  \node[text=NNHidden] at (68,-13) {$0.50$};

  \foreach \x/\word in {41/猫,55/追,69/狗}{
    \node at (\x,-36) {\word};
  }
  \foreach \r/\word in {0/猫,1/追,2/狗}{
    \node[anchor=east] at (28,-43-8*\r) {\word};
    \foreach \c in {0,1,2}{
      \ifnum\r=1
        \fill[NNHidden!28] (34+14*\c,-39-8*\r) rectangle ++(14,-8);
      \else
        \fill[NNHidden!12] (34+14*\c,-39-8*\r) rectangle ++(14,-8);
      \fi
      \draw[NNLine,line width=.5pt] (34+14*\c,-39-8*\r) rectangle ++(14,-8);
    }
  }
  \foreach \r/\a/\b/\c in {0/0.20/0.60/0.20,1/0.25/0.50/0.25,2/0.40/0.20/0.40}{
    \node at (41,-43-8*\r) {$\a$};
    \node at (55,-43-8*\r) {$\b$};
    \node at (69,-43-8*\r) {$\c$};
  }
  \draw[NNHidden,line width=1.1pt] (34,-47) rectangle (76,-55);
  \node at (55,-72)
    {$\bm z_2=0.25\bm v_1+0.50\bm v_2+0.25\bm v_3$};
\end{tikzpicture}
